\documentclass{article} % pas standalone
\usepackage{xparse}
\usepackage{tikz}
\include{pakage_tikz}
\usepackage{tikz}
\usepackage[europeanresistor]{circuitikz}
\usepgflibrary {shadings}
%\usetikzlibrary{calc}
\usetikzlibrary{hobby, decorations.markings, arrows.meta}
\usetikzlibrary{decorations.markings} 
\usetikzlibrary{calc,intersections,through,backgrounds}
\usepackage{amsmath}

%\usepackage{mathrsfs}

\usepackage{fancyhdr} % ajout du package fancyhdr

\pagestyle{fancy} % définition du style de page
\fancyhf{} % suppression des en-têtes et pieds de page actuels
\rfoot{\thepage} % ajout du numéro de page à droite
\fancyfoot[C]{\thepage}
\usepackage[margin=1cm]{geometry} % marges petites

% Commande qui crée une image sur une page entière
% --- Styles par défaut ---
% --- Styles par défaut ---
\tikzset{
  cadre/.style={draw=black!70, thick, rounded corners=3pt},
  ressort/.style={},
}

% --- Commande flexible : \TempImage[options]{width}{height}{content}
\newcommand{\TempImage}[4][]{%
  \begin{center}
    \begin{tikzpicture}[ressort, scale=1, transform shape]
      % --- Valeurs par défaut ---
      \def\bgcolor{gray!10} % couleur de fond par défaut
      \def\framecolor{black!60} % couleur du cadre
      \def\corner{3pt} % rayon des coins arrondis

      % --- Lecture des options facultatives ---
      \tikzset{tempimage/.cd, #1}

      % --- Cadre avec fond personnalisé ---
      \draw[
        fill=\bgcolor,
        draw=\framecolor,
        line width=1pt,
        rounded corners=\corner
      ] (0,0) rectangle (#2,#3);

      % --- Contenu centré ---
      \node[
        align=center,
        text width=#2-0.5cm,
        font=\small
      ] at (#2/2,#3/2) {#4};
    \end{tikzpicture}
  \end{center}
  \clearpage
}

% --- Définition des clés personnalisées ---
\tikzset{
  tempimage/.cd,
  bgcolor/.store in=\bgcolor,
  framecolor/.store in=\framecolor,
  corner/.store in=\corner,
}



\begin{document}

\tikzset{ressort/.style={thick,gray,smooth}}
\tikzstyle{cadre}=[rectangle,inner sep=5pt,inner ysep=5pt,draw,thick]


\TempImage{6cm}{4cm}{Equipe}

\TempImage{8cm}{3cm}{Piège harmonique anisotrope \\[0.3cm] Pièges optiques ou magnétiques \\[0.3cm] Confinement transversal fort, anisotrope \\[0.3cm] Gerbier (2004)}
\TempImage{5cm}{5cm}{Réseaux optiques 2D \\[0.3cm] Lattice de faisceaux croisés \\[0.3cm] Ensemble de tubes 1D \\[0.3cm] Paredes (2004), Kinoshita (2004)}
\TempImage{5cm}{5cm}{Pièges sur puce atomique \\[0.3cm] Fils micro-fabriqués, RF-dressing \\[0.3cm] Tube unique fortement confiné \\[0.3cm] van Amerongen (2008), Jacqmin (2011)}
\TempImage{5cm}{5cm}{Résonance de confinement (CIR) \\[0.3cm] Toute géométrie confinée \\[0.3cm] Modification du couplage $g_1D$\\[0.3cm] Olshanii (1998), Haller (2010)}

\TempImage[
    bgcolor=black!10!gray,
    framecolor=black!60!gray,
     corner=5pt
    ]{5cm}{5cm}{
    {\color{white!90!gray}
    \fontsize{5}{1}\selectfont
    One-dimensional quantum gases provide an ideal platform to study strong correlations and non-ergodic dynamics in out-of-equilibrium systems.\\[0.3cm]
Despite major progress, understanding the local rapidity distribution remains a key challenge.\\[0.3cm]
or \emph{Schematic view: from experimental system to theoretical description.}}}

\TempImage{8cm}{2cm}{Packet d'onde semiclassique en collisions}

\TempImage{5cm}{5cm}{graph (x,t) pour visualiser $\Delta$}


\TempImage{5cm}{7cm}{ Distribution $\rho(\theta)$\\[0.3cm]
Zoom sur $[\theta , \theta+ \delta \theta ]$ \\[0.3cm]
Bijection $\{I_a\} \leftrightarrow \{\theta_a\}$ \\[0.3cm] 
Veersion état fondamental }

\TempImage{8cm}{1cm}{Mer de Fermi + Exitations + nombre de Bethe $I_a$}

\TempImage{5cm}{7cm}{ Distribution $\rho(\theta)$\\[0.3cm]
Zoom sur $[\theta , \theta+ \delta \theta ]$ \\[0.3cm]
Bijection $\{I_a\} \leftrightarrow \{\theta_a\}$ \\[0.3cm] 
Veersion exité  }

\begin{tikzpicture}[>=stealth, node distance=2.3cm, thick, font=\small]

  % --- Styles ---
  \tikzstyle{exp} = [rectangle, rounded corners, draw=blue!60, fill=blue!10, minimum width=3cm, minimum height=1cm, align=center]
  \tikzstyle{theory} = [rectangle, rounded corners, draw=red!60, fill=red!10, minimum width=3cm, minimum height=1cm, align=center]
  \tikzstyle{process} = [ellipse, draw=orange!70!black, fill=orange!10, minimum width=2.2cm, align=center]

  % --- Nodes ---
  \node[exp] (exp) {1D Quantum Gas\\\small (Experiment)};
  \node[process, right of=exp] (measure) {Local\\Measurements};
  \node[theory, right of=measure] (rapidities) {Rapidity\\Distribution};
  \node[theory, right of=rapidities] (gge) {Generalized\\Gibbs Ensemble};

  % --- Arrows ---
  \draw[->, thick] (exp) -- (measure);
  \draw[->, thick] (measure) -- (rapidities);
  \draw[->, thick] (rapidities) -- (gge);

  % --- Optional labels ---
  \node[below=0.2cm of measure, align=center, font=\footnotesize] {Bridge between\\experiment and theory};
  \node[above=0.2cm of rapidities, align=center, font=\footnotesize] {Integrable\\description};

\end{tikzpicture}


\end{document}
